% Options for packages loaded elsewhere
% Options for packages loaded elsewhere
\PassOptionsToPackage{unicode}{hyperref}
\PassOptionsToPackage{hyphens}{url}
\PassOptionsToPackage{dvipsnames,svgnames,x11names}{xcolor}
%
\documentclass[
  spanish,
  a4paper,
  oneside]{scrbook}
\usepackage{xcolor}
\usepackage[top=25mm,bottom=25mm,left=30mm,right=20mm]{geometry}
\usepackage{amsmath,amssymb}
\setcounter{secnumdepth}{-\maxdimen} % remove section numbering
\usepackage{iftex}
\ifPDFTeX
  \usepackage[T1]{fontenc}
  \usepackage[utf8]{inputenc}
  \usepackage{textcomp} % provide euro and other symbols
\else % if luatex or xetex
  \usepackage{unicode-math} % this also loads fontspec
  \defaultfontfeatures{Scale=MatchLowercase}
  \defaultfontfeatures[\rmfamily]{Ligatures=TeX,Scale=1}
\fi
\usepackage{lmodern}
\ifPDFTeX\else
  % xetex/luatex font selection
\fi
% Use upquote if available, for straight quotes in verbatim environments
\IfFileExists{upquote.sty}{\usepackage{upquote}}{}
\IfFileExists{microtype.sty}{% use microtype if available
  \usepackage[]{microtype}
  \UseMicrotypeSet[protrusion]{basicmath} % disable protrusion for tt fonts
}{}
\makeatletter
\@ifundefined{KOMAClassName}{% if non-KOMA class
  \IfFileExists{parskip.sty}{%
    \usepackage{parskip}
  }{% else
    \setlength{\parindent}{0pt}
    \setlength{\parskip}{6pt plus 2pt minus 1pt}}
}{% if KOMA class
  \KOMAoptions{parskip=half}}
\makeatother
% Make \paragraph and \subparagraph free-standing
\makeatletter
\ifx\paragraph\undefined\else
  \let\oldparagraph\paragraph
  \renewcommand{\paragraph}{
    \@ifstar
      \xxxParagraphStar
      \xxxParagraphNoStar
  }
  \newcommand{\xxxParagraphStar}[1]{\oldparagraph*{#1}\mbox{}}
  \newcommand{\xxxParagraphNoStar}[1]{\oldparagraph{#1}\mbox{}}
\fi
\ifx\subparagraph\undefined\else
  \let\oldsubparagraph\subparagraph
  \renewcommand{\subparagraph}{
    \@ifstar
      \xxxSubParagraphStar
      \xxxSubParagraphNoStar
  }
  \newcommand{\xxxSubParagraphStar}[1]{\oldsubparagraph*{#1}\mbox{}}
  \newcommand{\xxxSubParagraphNoStar}[1]{\oldsubparagraph{#1}\mbox{}}
\fi
\makeatother


\usepackage{longtable,booktabs,array}
\newcounter{none} % for unnumbered tables
\usepackage{calc} % for calculating minipage widths
% Correct order of tables after \paragraph or \subparagraph
\usepackage{etoolbox}
\makeatletter
\patchcmd\longtable{\par}{\if@noskipsec\mbox{}\fi\par}{}{}
\makeatother
% Allow footnotes in longtable head/foot
\IfFileExists{footnotehyper.sty}{\usepackage{footnotehyper}}{\usepackage{footnote}}
\makesavenoteenv{longtable}
\usepackage{graphicx}
\makeatletter
\newsavebox\pandoc@box
\newcommand*\pandocbounded[1]{% scales image to fit in text height/width
  \sbox\pandoc@box{#1}%
  \Gscale@div\@tempa{\textheight}{\dimexpr\ht\pandoc@box+\dp\pandoc@box\relax}%
  \Gscale@div\@tempb{\linewidth}{\wd\pandoc@box}%
  \ifdim\@tempb\p@<\@tempa\p@\let\@tempa\@tempb\fi% select the smaller of both
  \ifdim\@tempa\p@<\p@\scalebox{\@tempa}{\usebox\pandoc@box}%
  \else\usebox{\pandoc@box}%
  \fi%
}
% Set default figure placement to htbp
\def\fps@figure{htbp}
\makeatother


% definitions for citeproc citations
\NewDocumentCommand\citeproctext{}{}
\NewDocumentCommand\citeproc{mm}{%
  \begingroup\def\citeproctext{#2}\cite{#1}\endgroup}
\makeatletter
 % allow citations to break across lines
 \let\@cite@ofmt\@firstofone
 % avoid brackets around text for \cite:
 \def\@biblabel#1{}
 \def\@cite#1#2{{#1\if@tempswa , #2\fi}}
\makeatother
\newlength{\cslhangindent}
\setlength{\cslhangindent}{1.5em}
\newlength{\csllabelwidth}
\setlength{\csllabelwidth}{3em}
\newenvironment{CSLReferences}[2] % #1 hanging-indent, #2 entry-spacing
 {\begin{list}{}{%
  \setlength{\itemindent}{0pt}
  \setlength{\leftmargin}{0pt}
  \setlength{\parsep}{0pt}
  % turn on hanging indent if param 1 is 1
  \ifodd #1
   \setlength{\leftmargin}{\cslhangindent}
   \setlength{\itemindent}{-1\cslhangindent}
  \fi
  % set entry spacing
  \setlength{\itemsep}{#2\baselineskip}}}
 {\end{list}}
\usepackage{calc}
\newcommand{\CSLBlock}[1]{\hfill\break\parbox[t]{\linewidth}{\strut\ignorespaces#1\strut}}
\newcommand{\CSLLeftMargin}[1]{\parbox[t]{\csllabelwidth}{\strut#1\strut}}
\newcommand{\CSLRightInline}[1]{\parbox[t]{\linewidth - \csllabelwidth}{\strut#1\strut}}
\newcommand{\CSLIndent}[1]{\hspace{\cslhangindent}#1}

\ifLuaTeX
\usepackage[bidi=basic,shorthands=off]{babel}
\else
\usepackage[bidi=default,shorthands=off]{babel}
\fi
\ifLuaTeX
  \usepackage{selnolig} % disable illegal ligatures
\fi


\setlength{\emergencystretch}{3em} % prevent overfull lines

\providecommand{\tightlist}{%
  \setlength{\itemsep}{0pt}\setlength{\parskip}{0pt}}



 


\AtBeginDocument{\hypersetup{linkcolor=black}}
% ---------- Paquetes tipograficos y microajustes ----------
\usepackage{microtype}        % mejor interletrado y justificado
\usepackage{csquotes}         % comillas tipograficas
\usepackage{iftex}
\ifPDFTeX\else
  % Seleccion de fuentes solo cuando XeTeX/LuaTeX esta disponible.
  \IfFontExistsTF{TeX Gyre Termes}{
    \setmainfont{TeX Gyre Termes}
  }{
    \IfFontExistsTF{Latin Modern Roman}{\setmainfont{Latin Modern Roman}}{}
  }
  \IfFontExistsTF{TeX Gyre Heros}{
    \setsansfont{TeX Gyre Heros}
  }{
    \IfFontExistsTF{Latin Modern Sans}{\setsansfont{Latin Modern Sans}}{}
  }
  \IfFontExistsTF{TeX Gyre Cursor}{
    \setmonofont{TeX Gyre Cursor}
  }{
    \IfFontExistsTF{Latin Modern Mono}{\setmonofont{Latin Modern Mono}}{}
  }
\fi
\usepackage{enumitem}         % listas compactas
\setlist{itemsep=.2em, topsep=.2em}

% ---------- KOMA-Script: estilo de titulos y espaciado ----------
\KOMAoptions{
  headings=big,
  parskip=half,
  fontsize=12pt,
  appendixprefix=true
}
\setlength{\parindent}{1.5em}
\setlength{\parskip}{0.6em}
\usepackage{indentfirst}
\usepackage{setspace}         % Control del interlineado del documento
\setstretch{1.5}

% ---------- Encabezados y pies (scrlayer-scrpage) ----------
\usepackage[automark,headsepline]{scrlayer-scrpage}
\clearpairofpagestyles
\automark[chapter]{chapter}
\ihead{\pagemark}
\ohead{\itshape\headmark}
\setheadsepline{0.4pt}
\renewcommand*{\chaptermarkformat}{}
\renewcommand*{\chapterpagestyle}{scrheadings}
\pagestyle{scrheadings}
\makeatletter
\let\ps@plain\ps@scrheadings
\makeatother
\cfoot{}

% ---------- Hipervinculos mas sobrios ----------
\usepackage{hyperref}
\hypersetup{
  colorlinks=true,
  linkcolor=black,
  citecolor=[rgb]{0.05,0.2,0.5},
  urlcolor=[rgb]{0.05,0.2,0.5},
  pdfauthor={\@author},
  pdftitle={\@title}
}

% ---------- Leyendas de figuras/tablas ----------
\usepackage[labelfont=bf,textfont=it]{caption}
\captionsetup{
  skip=8pt
}

% ---------- Tabla de contenidos ----------
\KOMAoptions{toc=graduated}
\RedeclareSectionCommand[tocnumwidth=3em]{chapter}
\RedeclareSectionCommand[tocindent=3.25em,tocnumwidth=2.8em]{section}
\RedeclareSectionCommand[tocindent=6.5em,tocnumwidth=2.5em]{subsection}
\makeatletter
\newcommand*{\tocdotfill}{\leavevmode\leaders\hbox to .6em{\hss.\hss}\hfill}
\makeatother
\RedeclareSectionCommand[toclinefill=\tocdotfill]{chapter}
\RedeclareSectionCommand[toclinefill=\tocdotfill]{section}
\RedeclareSectionCommand[toclinefill=\tocdotfill]{subsection}
\renewcommand*{\contentsname}{Tabla de contenido}
\setkomafont{chapterentry}{\normalfont}
\setkomafont{chapterentrypagenumber}{\normalfont}

% ---------- Viudas/Huerfanas y cortes de pagina ----------
\clubpenalty=10000
\widowpenalty=10000
\displaywidowpenalty=10000

% ---------- Entorno abstract para scrbook ----------
\providecommand{\abstractname}{Resumen}
\makeatletter
\@ifundefined{abstract}{
  \newenvironment{abstract}{
    \cleardoublepage
    \thispagestyle{plain}
    \null\vfill
    \begin{center}
      {\bfseries\Large \abstractname\par}
    \end{center}\vspace{1em}
    \begingroup
  }{
    \par\endgroup
    \vfill\null
    \cleardoublepage
  }
}{}
\makeatother

% ---------- Soporte de subtitulo desde YAML ----------
\makeatletter
\providecommand{\subtitle}[1]{\gdef\@subtitle{#1}}
\providecommand{\@subtitle}{}
\makeatother

\makeatletter
\newcommand{\facso@iflist}[4]{%
  \begingroup
    \def\addvspace##1{}%
    \def\numberline##1{##1}%
    \def\facso@target{#2}%
    \newcount\facso@entrycount
    \facso@entrycount=0\relax
    \long\def\contentsline##1##2##3{%
      \def\facso@this{##1}%
      \ifx\facso@this\facso@target
        \advance\facso@entrycount\@ne
      \fi
    }%
    \InputIfFileExists{\jobname.#1}{}{}%
  \endgroup
  \ifnum\facso@entrycount>0\relax
    #3%
  \else
    #4%
  \fi}
\renewcommand{\listoffigures}{%
  \facso@iflist{lof}{figure}{%
    \cleardoublepage
    \phantomsection
    \addchap*{Listado de Figuras}%
    \thispagestyle{scrheadings}%
    \markboth{Listado de Figuras}{Listado de Figuras}%
    \addcontentsline{toc}{chapter}{Listado de Figuras}%
    \@starttoc{lof}%
    \cleardoublepage
    \automark[chapter]{chapter}%
  }{}}
\renewcommand{\listoftables}{%
  \facso@iflist{lot}{table}{%
    \cleardoublepage
    \phantomsection
    \addchap*{Listado de Tablas}%
    \thispagestyle{scrheadings}%
    \markboth{Listado de Tablas}{Listado de Tablas}%
    \addcontentsline{toc}{chapter}{Listado de Tablas}%
    \@starttoc{lot}%
    \cleardoublepage
    \automark[chapter]{chapter}%
  }{}}
\makeatother

% --- Desactivar portada y abstract automaticos de Pandoc (PDF) ---
\AtBeginDocument{\let\maketitle\relax}
\renewenvironment{abstract}{}{}
\providecommand{\appendixname}{}
\providecommand{\appendixtocname}{}
\providecommand{\appendixpagename}{}
\renewcommand*{\appendixname}{Anexo}
\renewcommand*{\appendixtocname}{Anexos}
\renewcommand*{\appendixpagename}{Anexos}
\makeatletter
\@ifpackageloaded{bookmark}{}{\usepackage{bookmark}}
\makeatother
\makeatletter
\@ifpackageloaded{caption}{}{\usepackage{caption}}
\AtBeginDocument{%
\ifdefined\contentsname
  \renewcommand*\contentsname{Tabla de contenidos}
\else
  \newcommand\contentsname{Tabla de contenidos}
\fi
\ifdefined\listfigurename
  \renewcommand*\listfigurename{Lista de figuras}
\else
  \newcommand\listfigurename{Lista de figuras}
\fi
\ifdefined\listtablename
  \renewcommand*\listtablename{Lista de tablas}
\else
  \newcommand\listtablename{Lista de tablas}
\fi
\ifdefined\figurename
  \renewcommand*\figurename{Figura}
\else
  \newcommand\figurename{Figura}
\fi
\ifdefined\tablename
  \renewcommand*\tablename{Tabla}
\else
  \newcommand\tablename{Tabla}
\fi
}
\@ifpackageloaded{float}{}{\usepackage{float}}
\floatstyle{ruled}
\@ifundefined{c@chapter}{\newfloat{codelisting}{h}{lop}}{\newfloat{codelisting}{h}{lop}[chapter]}
\floatname{codelisting}{Listado}
\newcommand*\listoflistings{\listof{codelisting}{Listado de Listados}}
\makeatother
\makeatletter
\makeatother
\makeatletter
\@ifpackageloaded{caption}{}{\usepackage{caption}}
\@ifpackageloaded{subcaption}{}{\usepackage{subcaption}}
\makeatother
\usepackage{bookmark}
\IfFileExists{xurl.sty}{\usepackage{xurl}}{} % add URL line breaks if available
\urlstyle{same}
\hypersetup{
  pdftitle={¿Las mismas oportunidades para todos?},
  pdfauthor={Nicolás Outerbridge},
  pdflang={es},
  colorlinks=true,
  linkcolor={Maroon},
  filecolor={Maroon},
  citecolor={Blue},
  urlcolor={Blue},
  pdfcreator={LaTeX via pandoc}}


\title{¿Las mismas oportunidades para todos?}
\usepackage{etoolbox}
\makeatletter
\providecommand{\subtitle}[1]{% add subtitle to \maketitle
  \apptocmd{\@title}{\par {\large #1 \par}}{}{}
}
\makeatother
\subtitle{Creencias meritocráticas, experiencia escolar y percepción de
igualdad de oportunidades en estudiantes chilenos}
\author{Nicolás Outerbridge}
\date{22 de octubre de 2026}
\begin{document}
\frontmatter
\maketitle

% --- title-pdf.tex personalizado para portada formal ---
\makeatletter
\providecommand{\subtitle}[1]{\gdef\@subtitle{#1}}
\providecommand{\@subtitle}{}
\providecommand{\frontmattercontext}{}
\providecommand{\advisorname}{}
\providecommand{\advisorlabel}{Profesor guia:}
\providecommand{\frontmatterlocation}{}
\InputIfFileExists{includes/cover-config.tex}{}{}

\newcommand{\PrintTitle}{%
  {\sffamily\bfseries\fontsize{24pt}{28pt}\selectfont \@title\par}%
}
\newcommand{\PrintSubtitle}{%
  \begingroup
  \edef\temp{\detokenize{\@subtitle}}%
  \ifx\temp\empty\relax
    % sin subtitulo
  \else
    {\sffamily\bfseries\large \@subtitle\par}%
  \fi
  \endgroup
}
\newcommand{\PrintAuthor}{%
  {\Large\bfseries \@author\par}%
}
\newcommand{\PrintDate}{%
  {\small \@date\par}%
}
\makeatother

\begin{titlepage}
\thispagestyle{empty}
\begin{center}
\vspace*{10mm}

% Logo o imagen institucional
\includegraphics[width=0.35\textwidth]{assets/cover.png}\par
\vspace{12mm}

% Titulo y subtitulo
\PrintTitle
\vspace{6mm}
\PrintSubtitle

\vspace{22mm}
\begingroup
\edef\temp{\detokenize{\frontmattercontext}}%
\ifx\temp\empty\relax
  % sin contexto adicional
\else
  {\normalsize \frontmattercontext\par}
\fi
\endgroup

\vspace{18mm}
\PrintAuthor

\vspace{16mm}
\begingroup
\edef\temp{\detokenize{\advisorname}}%
\ifx\temp\empty\relax
  % sin tutor
\else
  \rule{0.45\textwidth}{0.4pt}\par
  {\small \advisorlabel\ \advisorname\par}
\fi
\endgroup

\vfill
\begingroup
\edef\temp{\detokenize{\frontmatterlocation}}%
\ifx\temp\empty\relax
  % sin ubicacion
\else
  {\small \frontmatterlocation\par}
\fi
\endgroup
\PrintDate

\end{center}
\end{titlepage}

% Preliminares (numeros romanos) y estilo simple
\frontmatter
\pagestyle{scrheadings}

\setcounter{tocdepth}{2} % (o 1/3 segun prefieras)
\tableofcontents
\listoffigures
\listoftables


\mainmatter
\bookmarksetup{startatroot}

\chapter*{Resumen}\label{resumen}
\addcontentsline{toc}{chapter}{Resumen}

\markboth{Resumen}{Resumen}

\small\textbf{Palabras clave:} igualdad de oportunidades, meritocracia,
privilegio, socialización escolar, educación superior, Chile \normalsize

La escuela promete que todos los estudiantes pueden llegar tan lejos
como su esfuerzo y su talento les permitan. Esa promesa descansa en un
supuesto anterior al mérito: que todos parten con las mismas
oportunidades. En Chile, sin embargo, el acceso a la educación superior
está fuertemente estratificado según el tipo de colegio de origen, y las
universidades más selectivas reclutan mayoritariamente a egresados de
colegios particulares pagados. Este estudio examina cómo perciben la
igualdad de oportunidades los estudiantes de básica y media, y qué
factores se asocian con esa percepción. Se analizan cuatro fuentes de
variación: las creencias meritocráticas de los estudiantes,
distinguiendo percepciones y preferencias sobre el mérito y el
privilegio; la experiencia de justicia en la propia escuela; la posición
en la jerarquía escolar y el horizonte educativo; y la exposición a
información sobre la desigualdad en el acceso a la universidad, mediante
un experimento de encuesta. Para ello se utilizan datos del panel
Educación y Meritocracia (EDUMER), aplicado a estudiantes de básica y
media de nueve escuelas chilenas entre 2023 y 2024, cuya primera ola
incluye un experimento con asignación aleatoria de información.

\bookmarksetup{startatroot}

\chapter*{Prefacio}\label{prefacio}
\addcontentsline{toc}{chapter}{Prefacio}

\markboth{Prefacio}{Prefacio}

\mainmatter

\bookmarksetup{startatroot}

\chapter{Introducción}\label{introducciuxf3n}

Las instituciones educativas tienen la función explícita de transmitir
conocimientos y habilidades que preparan a los estudiantes para sus
futuras posiciones en la sociedad. Asimismo, transmiten normas, valores
y creencias que respaldan el logro educativo
(\citeproc{ref-mijs2016}{Mijs, 2016b}). Entre los principios centrales
destaca la promesa de que la escuela ofrece a todos la misma oportunidad
de demostrar su valía, y que las diferencias en los resultados reflejan
diferencias de esfuerzo y talento, no de origen
(\citeproc{ref-batruch_belief_2023}{Batruch et~al., 2023};
\citeproc{ref-bourdieu_reproduction_1990}{Bourdieu \& Passeron, 1990}).

Esta promesa comprende dos elementos diferentes. El primero es la
\emph{meritocracia}, definida como la regla que distribuye recompensas
según esfuerzo y talento (\citeproc{ref-young1958}{Young, 1958}). El
segundo, que podría suponer un supuesto anterior, corresponde a la
\emph{igualdad de oportunidades}, entendida como la condición en la que
todos compiten desde puntos de partida comparables, de modo que el
resultado dependa de acciones individuales y no del origen social. Solo
en un contexto de igualdad de oportunidades la meritocracia puede
generar movilidad social
(\citeproc{ref-castillo_meritocracia_2019}{Castillo et~al., 2019}). Una
sociedad puede recompensar el esfuerzo y, simultáneamente, distribuir
desigualmente las oportunidades para alcanzarlo. Según Mijs
(\citeproc{ref-mijs2016}{2016b}), las oportunidades para demostrar
mérito dependen de factores externos a la meritocracia, lo que implica
que esta no solo es una promesa incumplida, sino también incumplible.

El caso chileno ilustra esta tensión. El sistema escolar se caracteriza
por una alta estratificación y la noción de mérito ocupa un lugar
central en el discurso educativo
(\citeproc{ref-castillo_socialization_2024}{Castillo, Salgado, et~al.,
2024}). De manera paralela, la educación superior experimentó una rápida
masificación bajo un sistema jerarquizado, en el que pocas instituciones
forman a las élites y otras absorben a la mayoría
(\citeproc{ref-brunner_lucha_2012}{Brunner, 2012};
\citeproc{ref-brunner_dinamicas_2016}{Brunner \& Ganga, 2016}). Las
universidades más selectivas reclutan principalmente a egresados de
colegios particulares pagados, que representan una fracción menor de la
matrícula escolar total (\citeproc{ref-brunner_lucha_2012}{Brunner,
2012}; \citeproc{ref-munoz_taxonomia_2013}{Muñoz \& Blanco, 2013}). Para
los estudiantes de educación básica o media, la cuestión de si todos
tienen las mismas oportunidades no es abstracta, sino que se concreta en
el horizonte de acceso a distintas instituciones al finalizar la
escuela.

La investigación sobre creencias distributivas en la población escolar
se ha enfocado principalmente en la meritocracia, es decir, en si los
estudiantes creen que el esfuerzo es recompensado y si consideran que
debería serlo (\citeproc{ref-castillo_socialization_2024}{Castillo,
Salgado, et~al., 2024}; \citeproc{ref-laffert_changes}{Laffert et~al.,
s.f.}). Esta línea de investigación ha mostrado que las creencias
meritocráticas son multidimensionales: los estudiantes distinguen entre
lo que perciben y lo que prefieren, y entre los principios del mérito y
del privilegio, de modo que pueden adherir al mérito como ideal y, al
mismo tiempo, reconocer el peso de la riqueza familiar y los contactos
(\citeproc{ref-laffert_changes}{Laffert et~al., s.f.}). La percepción de
igualdad de oportunidades, en cambio, ha recibido menor atención y, con
frecuencia, se ha tratado como equivalente a la creencia meritocrática
(\citeproc{ref-laffert_changes}{Laffert et~al., s.f.}). Además, la
evidencia sobre cómo los estudiantes perciben la desigualdad en el
acceso a la educación superior proviene mayoritariamente de quienes ya
han ingresado a ella, especialmente estudiantes de universidades de
élite (\citeproc{ref-allouch_ethique_2016}{Allouch et~al., 2016};
\citeproc{ref-binder_symbolically_2019}{Binder \& Abel, 2019};
\citeproc{ref-villalobos_percepciones_2022}{Villalobos et~al., 2022}).
Existe escaso conocimiento sobre la formación de esta percepción en
etapas previas, cuando los estudiantes aún se encuentran en la escuela y
su trayectoria futura permanece abierta.

El presente estudio aborda esta brecha mediante la siguiente pregunta:
¿cómo perciben la igualdad de oportunidades los estudiantes escolares
chilenos y en qué medida dicha percepción se asocia con sus creencias
sobre el mérito y el privilegio, su experiencia de justicia en la
escuela, su posición en la jerarquía escolar y su horizonte educativo?
Asimismo, se examina si proporcionar información sobre la desigualdad en
el acceso a la universidad modifica su percepción sobre la igualdad de
oportunidades. Para responder a estas cuestiones, se utilizan datos del
panel Educación y Meritocracia (EDUMER), cuya primera ola incluye un
experimento en el que la mitad de los estudiantes recibió información
sobre diferencias de acceso a la universidad según el tipo de colegio.

\bookmarksetup{startatroot}

\chapter{Antecedentes conceptuales y
empíricos}\label{antecedentes-conceptuales-y-empuxedricos}

\section{Igualdad de oportunidades y
meritocracia}\label{igualdad-de-oportunidades-y-meritocracia}

La investigación sobre creencias meritocráticas ha distinguido, en
primer lugar, entre una dimensión normativa, la \emph{preferencia} por
que el esfuerzo y el talento sean la base de la distribución de
recompensas, y una dimensión descriptiva, la \emph{percepción} de que la
sociedad efectivamente los recompensa
(\citeproc{ref-castillo_meritocracia_2019}{Castillo et~al., 2019};
\citeproc{ref-duru-bellat_whos_2012}{Duru-Bellat \& Tenret, 2012}). La
percepción se refiere a lo que es observado (``lo que es'') y la
preferencia a lo que las personas quisieran que ocurriera (``lo que
debería ser'') (\citeproc{ref-castillo_meritocracia_2019}{Castillo
et~al., 2019}). Ambas dimensiones constituyen constructos diferentes: en
población adulta chilena, su correlación es baja y no significativa, y
la preferencia es más alta, más consensual y resistente, de modo que
puede mantenerse aunque se perciba que la meritocracia no opera
(\citeproc{ref-castillo_meritocracia_2019}{Castillo et~al., 2019}).

Castillo et al. (\citeproc{ref-castillo_multidimensional_2023}{2023})
agregan una segunda distinción, entre los elementos meritocráticos
(esfuerzo y talento) y los elementos basados en el privilegio (riqueza
familiar y contactos), que se tratan como componentes distintos y no
como polos opuestos de un mismo continuo
(\citeproc{ref-laffert_changes}{Laffert et~al., s.f.}). De la
combinación de ambas distinciones resultan cuatro dimensiones: la
\emph{percepción de meritocracia} (si el esfuerzo y la habilidad son
recompensados), la \emph{percepción de privilegio} (si el éxito se
vincula a los contactos y la riqueza familiar), la \emph{preferencia
meritocrática} (si quienes se esfuerzan más o tienen más talento
deberían obtener mayores recompensas) y la \emph{aceptación del
privilegio} (si es aceptable que quienes tienen mejores contactos o
padres ricos obtengan mayor éxito)
(\citeproc{ref-laffert_changes}{Laffert et~al., s.f.}). Esta estructura
de cuatro factores se ha validado en población escolar chilena con los
mismos datos que utiliza este estudio, y los resultados muestran que la
adhesión al mérito coexiste con el reconocimiento del privilegio, en
lugar de desplazarlo (\citeproc{ref-laffert_changes}{Laffert et~al.,
s.f.}).

Los niveles de cada dimensión ilustran esta coexistencia. En la primera
ola del panel EDUMER, la mayoría de los estudiantes percibe que en Chile
les va mejor a quienes tienen buenos contactos (77\%) o padres ricos
(71\%), y en menor medida que se recompensa el talento (74\%) o el
esfuerzo (67\%). Al mismo tiempo, el 87\% considera que el esfuerzo
debería ser la base de las recompensas, mientras que la aceptación de
que a quienes tienen padres ricos les vaya mejor se divide casi por
mitades y la de los buenos contactos es mayoritaria (61\%)
(\citeproc{ref-laffert_changes}{Laffert et~al., s.f.}). Las cuatro
dimensiones se relacionan entre sí de manera diferenciada: la percepción
de meritocracia y la de privilegio se correlacionan negativamente,
mientras que la percepción de privilegio se asocia positivamente con la
preferencia meritocrática, lo que se ha interpretado como una respuesta
normativo-reactiva ante la injusticia percibida
(\citeproc{ref-laffert_changes}{Laffert et~al., s.f.}).

La percepción de igualdad de oportunidades es también una creencia
descriptiva, pero se refiere a otro momento del proceso distributivo.
Mientras la percepción de meritocracia evalúa la regla que conecta el
desempeño con la recompensa, la percepción de igualdad de oportunidades
evalúa las condiciones de partida: si el origen familiar, el colegio o
los recursos del hogar determinan quién puede competir y con qué
ventajas (\citeproc{ref-mijs2016}{Mijs, 2016b}). Por ello, ambas
creencias pueden disociarse. Un estudiante puede creer que en Chile el
esfuerzo es recompensado y, al mismo tiempo, reconocer que no todos
parten desde el mismo lugar.

Parte de la literatura ha operacionalizado la creencia meritocrática
mediante el apoyo genérico a la igualdad de oportunidades
(\citeproc{ref-batruch_belief_2023}{Batruch et~al., 2023};
\citeproc{ref-darnon_where_2018}{Darnon et~al., 2018}), lo que dificulta
distinguir ambas creencias (\citeproc{ref-laffert_changes}{Laffert
et~al., s.f.}). Incluso las mediciones de la preferencia meritocrática
han incorporado la igualdad de oportunidades como condición, por
ejemplo, al preguntar si está bien que algunas personas tengan más
dinero ``pero sólo si hay igualdad de oportunidades''
(\citeproc{ref-castillo_meritocracia_2019}{Castillo et~al., 2019}).
Separar ambas creencias tiene consecuencias sustantivas. En la
investigación sobre percepciones de desigualdad, distintos indicadores
perceptuales se asocian débilmente y capturan aspectos diferentes del
fenómeno, con determinantes distintos
(\citeproc{ref-castillo_percepcion_2012}{Castillo et~al., 2012}). Este
estudio se centra en la percepción de igualdad de oportunidades y trata
las cuatro dimensiones de la creencia meritocrática como creencias
relacionadas pero distintas, que pueden contribuir a explicarla.

\section{Creencias meritocráticas y percepción de igualdad de
oportunidades}\label{creencias-meritocruxe1ticas-y-percepciuxf3n-de-igualdad-de-oportunidades}

Las cuatro dimensiones permiten derivar expectativas distintas sobre la
percepción de igualdad de oportunidades. La meritocracia opera a la vez
como un juicio cognitivo sobre cómo funciona la desigualdad y como un
lente moral para evaluar si los resultados desiguales son merecidos
(\citeproc{ref-laffert_changes}{Laffert et~al., s.f.}). Por ello, sus
dimensiones no tienen por qué relacionarse en la misma dirección con la
evaluación de las oportunidades.

La percepción de meritocracia cumpliría una función de legitimación.
Quienes perciben que la meritocracia funciona adecuadamente perciben
menos desigualdad económica, lo que Castillo et al.
(\citeproc{ref-castillo_meritocracia_2019}{2019}) interpretan como una
legitimación de las diferencias de ingreso, que se deberían a razones
más bien individuales que sistémicas. Una percepción alta de
meritocracia es más compatible con la afirmación del orden existente,
mientras que una percepción baja puede reflejar insatisfacción y una
demanda de mayor justicia (\citeproc{ref-laffert_changes}{Laffert
et~al., s.f.}). Si las recompensas se perciben como resultado del
esfuerzo y el talento, es esperable que también se perciba que las
condiciones para competir están abiertas a todos.

La percepción de privilegio apunta en la dirección contraria. Percibir
que a quienes tienen padres ricos o buenos contactos les va mejor
equivale a reconocer que el origen familiar y las redes condicionan los
resultados, es decir, que las oportunidades para demostrar mérito
dependen de factores externos a él (\citeproc{ref-mijs2016}{Mijs,
2016b}). La percepción de privilegio debería, por tanto, asociarse con
una menor percepción de igualdad de oportunidades.

La preferencia meritocrática no implica una evaluación positiva de la
realidad. En población adulta, una mayor preferencia por la meritocracia
se asocia con una mayor percepción de desigualdad, lo que sugiere que la
adhesión normativa al mérito establece un estándar más exigente frente a
la realidad y una demanda por reducir las diferencias que se basan en
criterios distintos al mérito
(\citeproc{ref-castillo_meritocracia_2019}{Castillo et~al., 2019}). En
población escolar, la preferencia meritocrática se asocia con una mayor
percepción de privilegio (\citeproc{ref-laffert_changes}{Laffert et~al.,
s.f.}). Quienes prefieren con más fuerza que el esfuerzo y el talento
determinen las recompensas tenderían, así, a evaluar más críticamente si
todos parten desde el mismo lugar.

La aceptación del privilegio es la dimensión menos estudiada. Su
relación con la percepción de privilegio es positiva, lo que indica que
reconocer mecanismos no meritocráticos se vincula con aceptar en mayor
medida su legitimidad (\citeproc{ref-laffert_changes}{Laffert et~al.,
s.f.}). La preferencia por la meritocracia ha sido descrita como un mito
legitimador, que provee justificaciones morales y racionales para la
desigualdad social (\citeproc{ref-castillo_meritocracia_2019}{Castillo
et~al., 2019}). Si se extiende este argumento a la aceptación del
privilegio, quienes consideran legítimo que el origen familiar y los
contactos favorezcan a algunos tendrían menos razones para interpretar
esas ventajas como una desigualdad de oportunidades. Al mismo tiempo,
como esta aceptación puede operar como un juicio normativo separado de
la evaluación del mérito (\citeproc{ref-laffert_changes}{Laffert et~al.,
s.f.}), esta expectativa tiene un carácter más exploratorio.

\section{La escuela como espacio de socialización de las
oportunidades}\label{la-escuela-como-espacio-de-socializaciuxf3n-de-las-oportunidades}

La escuela es la institución que de manera más explícita promete
igualdad de condiciones: todos reciben el mismo currículo, son evaluados
con las mismas pruebas y compiten por las mismas notas. Su estructura,
su discurso y sus prácticas de evaluación contribuyen a presentar el
éxito académico como resultado del talento y el esfuerzo individual, y a
transformar diferencias de origen en diferencias de desempeño que
aparecen como individuales
(\citeproc{ref-bourdieu_reproduction_1990}{Bourdieu \& Passeron, 1990};
\citeproc{ref-goudeau_hidden_2017}{Goudeau \& Croizet, 2017}).

Esta promesa no se cumple. Las oportunidades para demostrar mérito
dependen del origen social (\citeproc{ref-mijs2016}{Mijs, 2016b}), y las
prácticas de aula convierten ventajas de clase en diferencias de
rendimiento (\citeproc{ref-goudeau_hidden_2017}{Goudeau \& Croizet,
2017}). La experiencia de justicia escolar tampoco es homogénea: en
Israel, cerca de la mitad de los estudiantes de secundaria se considera
perjudicada en sus calificaciones (\citeproc{ref-resh2010}{Resh \&
Sabbagh, 2010}). La forma en que la escuela organiza a los estudiantes
también moldea sus explicaciones: en sistemas que separan a los
estudiantes según rendimiento, estos tienden a atribuir su fracaso a sí
mismos (\citeproc{ref-mijs_stratified_2016}{Mijs, 2016a}).

Lo que ocurre en la escuela parece extenderse a la evaluación de la
sociedad. Creer que la escuela es meritocrática se asocia con percibir
como menos injusta la desigualdad social y con un menor apoyo a la
redistribución, y esta relación se mantiene al controlar por la creencia
en la meritocracia general (\citeproc{ref-batruch_belief_2023}{Batruch
et~al., 2023}). Batruch et al.~interpretan que creer que la escuela es
justa implica creer que también lo son las diferencias que produce la
jerarquía educativa. Si la escuela es la primera institución en que los
estudiantes experimentan una distribución sistemática de recompensas, la
evaluación de que esa distribución es justa podría extenderse a la
percepción de que la sociedad ofrece a todos las mismas oportunidades.

\section{La educación superior como horizonte de
movilidad}\label{la-educaciuxf3n-superior-como-horizonte-de-movilidad}

Para los estudiantes escolares, la desigualdad de oportunidades se
materializa sobre todo en el acceso a la educación superior. Chile es un
caso paradigmático de masificación terciaria en América Latina
(\citeproc{ref-brunner_dinamicas_2016}{Brunner \& Ganga, 2016}): la
cobertura de pregrado pasó de cerca de 8\% en 1980 a 53\% en 2015, en un
crecimiento absorbido principalmente por instituciones privadas y con
escasa regulación estatal
(\citeproc{ref-villalobos_percepciones_2022}{Villalobos et~al., 2022}).
La masificación, sin embargo, no eliminó la estratificación del acceso.
Según la encuesta CASEN 2009, la tasa de participación en educación
terciaria variaba entre 22,7\% en el quintil de menores ingresos y
67,8\% en el de mayores ingresos
(\citeproc{ref-brunner_lucha_2012}{Brunner, 2012}).

Más importante aún, el sistema se organizó como una pirámide jerárquica
de selectividad y reputación. Brunner
(\citeproc{ref-brunner_lucha_2012}{2012}) describe el ingreso a las
universidades chilenas como un proceso de doble selección, académica y
socioeconómica, donde las instituciones más prestigiosas eligen a los
estudiantes con mejores puntajes y las demás se ordenan en cascada hasta
la base, donde no existe selectividad. En la cima, la composición
escolar del alumnado revela el origen familiar: dos tercios de quienes
ingresaban a la Pontificia Universidad Católica provenían de colegios
particulares pagados, que reunían apenas el 6,9\% de la matrícula
nacional de enseñanza media, mientras solo un 13\% provenía de liceos
municipales. Brunner resume esta dinámica señalando que las
universidades donde se forman las élites reproducen ``bajo un velo
meritocrático el privilegio de los herederos'', al tiempo que ofrecen un
canal de ascenso a quienes superan las barreras selectivas.

La taxonomía de Muñoz y Blanco
(\citeproc{ref-munoz_taxonomia_2013}{2013}), construida mediante
análisis factorial sobre variables de operación de las universidades,
confirma que el origen escolar del alumnado es uno de los ejes que
ordenan el sistema. De los cinco tipos que identifican, dos se definen
precisamente por la composición escolar de su matrícula: las
universidades \emph{elitistas}, con alta selectividad y una proporción
de egresados de colegios particulares pagados que supera el 70\% en las
``altamente elitistas'' (Universidad de los Andes, Adolfo Ibáñez y del
Desarrollo), y las universidades \emph{no elitistas}, con selectividad
menor al 2\% y alumnado proveniente principalmente de establecimientos
municipales. Entre ambos extremos se ubican las universidades de
investigación, las de acreditación y las masivas, estas últimas con más
de la mitad de su matrícula proveniente de colegios particulares
subvencionados.

En consecuencia, la pregunta por la igualdad de oportunidades adquiere
en Chile un contenido institucional preciso: el tipo de colegio anticipa
el tipo de institución a la que se puede aspirar. La experiencia
internacional muestra que esta jerarquía no desaparece con la expansión
del acceso. En el Reino Unido, la masificación aumentó la competencia
entre instituciones y postulantes, y los estudiantes de clase
trabajadora siguen concentrándose en instituciones de menor prestigio,
cuyos títulos tienen menor valor en el mercado laboral
(\citeproc{ref-bathmaker_degree_2016}{Bathmaker et~al., 2016}). La
desigualdad se prolonga después de la universidad: las firmas
profesionales de élite privilegian a candidatos con formas estrechas de
capital cultural, aun cuando declaran reclutar al ``mejor talento''
(\citeproc{ref-ashley_understanding_2017}{Ashley \& Empson, 2017}).

\section{Posición en la jerarquía escolar y horizonte
educativo}\label{posiciuxf3n-en-la-jerarquuxeda-escolar-y-horizonte-educativo}

La evidencia disponible sobre cómo los estudiantes perciben la
desigualdad en el acceso a la educación superior proviene de quienes ya
están dentro. En Chile, Villalobos et al.
(\citeproc{ref-villalobos_percepciones_2022}{2022}) encuestaron a 1.773
estudiantes de ocho carreras de universidades de élite. El principal eje
de diferenciación de sus percepciones no fue el nivel socioeconómico,
sino el tipo de universidad: los estudiantes de universidades
tradicionales eran más escépticos de la meritocracia y de las
explicaciones individualistas de la desigualdad que sus pares de
universidades privadas elitistas. Los autores observan además
percepciones ``ortogonales'': los estudiantes explican la desigualdad
general por factores estructurales, pero atribuyen su propio ingreso a
sus cualidades personales.

Estudios en otros países apuntan en la misma dirección. Los estudiantes
de Sciences Po y Oxford reconocen que viven en un mundo desigual y que
sus trayectorias están condicionadas por esa desigualdad, pero al mismo
tiempo creen firmemente que merecen el lugar que ocupan, y recurren a la
idea de responsabilidad para justificar su privilegio
(\citeproc{ref-allouch_ethique_2016}{Allouch et~al., 2016}). Los
estudiantes de Harvard y Stanford, por su parte, construyen fronteras
simbólicas entre universidades de élite que justifican su propia
posición (\citeproc{ref-binder_symbolically_2019}{Binder \& Abel,
2019}). Quienes acceden a las posiciones aventajadas tienden, así, a
interpretar su trayectoria como merecida.

En la escuela, la posición de cada estudiante se expresa ante todo en
sus calificaciones, que operan como marcadores institucionales de valor
y hacen visible el mérito de manera inmediata
(\citeproc{ref-laffert_changes}{Laffert et~al., s.f.}). Al presentar los
resultados como ganados, los ideales meritocráticos alientan a los
``ganadores'' a ver su posición como merecida
(\citeproc{ref-laffert_changes}{Laffert et~al., s.f.}). Desde una
perspectiva de \emph{interés propio}, quienes ocupan posiciones
ventajosas o son reconocidos por el sistema escolar tendrían razones
para percibir que las oportunidades están abiertas a todos, porque esa
percepción legitima su lugar (\citeproc{ref-baer_education_1982}{Baer \&
Lambert, 1982}; \citeproc{ref-dickinson_adolescent_1990}{Dickinson,
1990}; \citeproc{ref-duru-bellat_whos_2012}{Duru-Bellat \& Tenret,
2012}). Del mismo modo, quienes se proyectan en la universidad ocupan
una posición anticipada de ganadores del sistema.

\section{Trayectoria escolar: socialización e
ilustración}\label{trayectoria-escolar-socializaciuxf3n-e-ilustraciuxf3n}

Una pregunta distinta es cómo cambia la percepción de igualdad de
oportunidades a medida que los estudiantes avanzan en la escuela y se
acercan a la transición a la educación superior. La hipótesis de
\emph{socialización} plantea que la educación transmite la ideología
dominante según el tiempo que se pasa en ella; la de \emph{ilustración},
que fomenta una mirada crítica sobre la desigualdad
(\citeproc{ref-baer_education_1982}{Baer \& Lambert, 1982}). Los
estudios sobre socialización universitaria muestran que los contextos
educativos pueden moldear actitudes sociopolíticas de forma duradera
(\citeproc{ref-alwin_political_1991}{Alwin et~al., 1991}), y que la
exposición a la institución modifica las percepciones: entre estudiantes
de élite chilenos, quienes llevan más tiempo en su carrera atribuyen
menos su ingreso a sus atributos personales
(\citeproc{ref-villalobos_percepciones_2022}{Villalobos et~al., 2022}).

En la etapa escolar, la evidencia es escasa y ambigua. Entre escolares
chilenos, los estudiantes de media usan la escala de percepción de
meritocracia de manera distinta a los de básica al evaluar si el
esfuerzo es recompensado, lo que se interpreta como un marco más crítico
entre los mayores (\citeproc{ref-laffert_changes}{Laffert et~al.,
s.f.}). Es plausible que, a medida que la selección universitaria se
vuelve más próxima y visible, los estudiantes perciban con mayor
claridad que no todos parten desde el mismo lugar.

\section{Hipótesis}\label{hipuxf3tesis}

A partir de los antecedentes se plantean seis hipótesis, resumidas en la
Tabla~\ref{tbl-hipotesis}.

\textbf{Transferencia.} Si la escuela es la primera institución en que
los estudiantes experimentan una distribución de recompensas que se
presenta como igualitaria y basada en el mérito, la evaluación de que
esa distribución es justa debería extenderse a la percepción de que la
sociedad ofrece las mismas oportunidades a todos
(\citeproc{ref-batruch_belief_2023}{Batruch et~al., 2023};
\citeproc{ref-mijs_stratified_2016}{Mijs, 2016a}).

\begin{quote}
\textbf{H1:} A mayor percepción de meritocracia en la escuela, mayor
percepción de igualdad de oportunidades en la sociedad.
\end{quote}

\textbf{Posición en la jerarquía escolar.} Quienes son reconocidos por
el sistema escolar tendrían razones para percibir que las oportunidades
están abiertas a todos, porque esa percepción legitima su lugar
(\citeproc{ref-baer_education_1982}{Baer \& Lambert, 1982};
\citeproc{ref-duru-bellat_whos_2012}{Duru-Bellat \& Tenret, 2012};
\citeproc{ref-laffert_changes}{Laffert et~al., s.f.}).

\begin{quote}
\textbf{H2:} Los estudiantes con mejor rendimiento académico perciben
más igualdad de oportunidades.
\end{quote}

\textbf{Horizonte educativo.} Quienes se proyectan en la educación
superior, y en particular en la universidad, ocupan una posición
anticipada de ``ganadores'' del sistema. Como los estudiantes que ya
accedieron a instituciones selectivas, tenderían a interpretar su
trayectoria esperada como resultado de su esfuerzo y a percibir que las
oportunidades están abiertas
(\citeproc{ref-allouch_ethique_2016}{Allouch et~al., 2016};
\citeproc{ref-villalobos_percepciones_2022}{Villalobos et~al., 2022}).

\begin{quote}
\textbf{H3:} Los estudiantes que esperan completar una carrera
universitaria perciben más igualdad de oportunidades que quienes esperan
completar solo la educación media o una carrera técnica.
\end{quote}

\textbf{Información sobre el acceso a la universidad.} Si la percepción
de igualdad de oportunidades descansa en parte en el desconocimiento de
la estratificación del acceso a la educación superior, hacer visible esa
estratificación debería reducirla
(\citeproc{ref-brunner_lucha_2012}{Brunner, 2012};
\citeproc{ref-munoz_taxonomia_2013}{Muñoz \& Blanco, 2013}). Mientras H1
supone que la percepción se forma a partir de la experiencia en la
escuela, el experimento permite evaluar si la información sobre la
desigualdad del acceso basta para modificarla. La evidencia previa con
este experimento muestra efectos pequeños sobre la percepción de
meritocracia y dificultades para comprender la información
(\citeproc{ref-castillo_experimento_2024}{Castillo, Carrasco, et~al.,
2024}).

\begin{quote}
\textbf{H4:} Los estudiantes expuestos a información sobre las
diferencias de acceso a la universidad según tipo de colegio
(tratamiento) perciben menos igualdad de oportunidades que el grupo
control.
\end{quote}

\textbf{Socialización e ilustración.} La evidencia no permite anticipar
si el avance en la escolaridad refuerza o cuestiona la percepción de
igualdad de oportunidades (\citeproc{ref-baer_education_1982}{Baer \&
Lambert, 1982}; \citeproc{ref-duru-bellat_whos_2012}{Duru-Bellat \&
Tenret, 2012}; \citeproc{ref-laffert_changes}{Laffert et~al., s.f.}),
por lo que se plantea una hipótesis bilateral.

\begin{quote}
\textbf{H5:} La percepción de igualdad de oportunidades cambia con el
avance en la trayectoria escolar. Si predomina la socialización,
aumenta; si predomina la ilustración, disminuye.
\end{quote}

Para distinguir entre ambas explicaciones se examinará
complementariamente el cambio en la percepción de privilegio. La
ilustración implica que, junto con la menor percepción de igualdad de
oportunidades, aumenta la percepción de que el éxito depende de tener
padres ricos o buenos contactos.

\textbf{Creencias meritocráticas.} Se espera que las cuatro dimensiones
de la creencia meritocrática se relacionen de manera diferenciada con la
percepción de igualdad de oportunidades: la percepción de meritocracia
legitimaría las diferencias existentes, la percepción de privilegio
reconoce el peso del origen, la preferencia meritocrática establecería
un estándar más exigente frente a la realidad y la aceptación del
privilegio legitimaría las ventajas de origen
(\citeproc{ref-castillo_meritocracia_2019}{Castillo et~al., 2019};
\citeproc{ref-laffert_changes}{Laffert et~al., s.f.};
\citeproc{ref-mijs2016}{Mijs, 2016b}). H6d tiene un carácter
exploratorio.

\begin{quote}
\textbf{H6a:} A mayor percepción de meritocracia, mayor percepción de
igualdad de oportunidades.
\end{quote}

\begin{quote}
\textbf{H6b:} A mayor percepción de privilegio, menor percepción de
igualdad de oportunidades.
\end{quote}

\begin{quote}
\textbf{H6c:} A mayor preferencia meritocrática, menor percepción de
igualdad de oportunidades.
\end{quote}

\begin{quote}
\textbf{H6d:} A mayor aceptación del privilegio, mayor percepción de
igualdad de oportunidades.
\end{quote}

\begin{longtable}[]{@{}
  >{\raggedright\arraybackslash}p{(\linewidth - 6\tabcolsep) * \real{0.2500}}
  >{\raggedright\arraybackslash}p{(\linewidth - 6\tabcolsep) * \real{0.2500}}
  >{\raggedright\arraybackslash}p{(\linewidth - 6\tabcolsep) * \real{0.2500}}
  >{\raggedright\arraybackslash}p{(\linewidth - 6\tabcolsep) * \real{0.2500}}@{}}
\caption{Resumen de hipótesis}\label{tbl-hipotesis}\tabularnewline
\toprule\noalign{}
\begin{minipage}[b]{\linewidth}\raggedright
Hipótesis
\end{minipage} & \begin{minipage}[b]{\linewidth}\raggedright
Mecanismo
\end{minipage} & \begin{minipage}[b]{\linewidth}\raggedright
Predictor
\end{minipage} & \begin{minipage}[b]{\linewidth}\raggedright
Efecto esperado
\end{minipage} \\
\midrule\noalign{}
\endfirsthead
\toprule\noalign{}
\begin{minipage}[b]{\linewidth}\raggedright
Hipótesis
\end{minipage} & \begin{minipage}[b]{\linewidth}\raggedright
Mecanismo
\end{minipage} & \begin{minipage}[b]{\linewidth}\raggedright
Predictor
\end{minipage} & \begin{minipage}[b]{\linewidth}\raggedright
Efecto esperado
\end{minipage} \\
\midrule\noalign{}
\endhead
\bottomrule\noalign{}
\endlastfoot
H1 & Transferencia & Meritocracia escolar percibida & Positivo \\
H2 & Interés propio & Rendimiento académico & Positivo \\
H3 & Horizonte educativo & Expectativa de completar la universidad &
Positivo \\
H4 & Información sobre el acceso & Asignación al tratamiento &
Negativo \\
H5 & Socialización / ilustración & Ola (panel) & Positivo
(socialización) / negativo (ilustración) \\
H6a & Legitimación & Percepción de meritocracia & Positivo \\
H6b & Reconocimiento del origen & Percepción de privilegio & Negativo \\
H6c & Estándar normativo & Preferencia meritocrática & Negativo \\
H6d & Legitimación del privilegio & Aceptación del privilegio &
Positivo \\
\end{longtable}

\bookmarksetup{startatroot}

\chapter{Metodología}\label{metodologuxeda}

\section{Datos}\label{datos}

Se utilizan datos de la Encuesta de Percepciones de Meritocracia y
Formación Ciudadana, levantada en el marco del panel Educación y
Meritocracia (EDUMER), proyecto FONDECYT N.º 1210847
(\citeproc{ref-edumer_manual}{Proyecto Educación y Meritocracia, s.f.}).
Su objetivo es examinar las creencias, actitudes y comportamientos de
estudiantes en edad escolar respecto de la meritocracia, la desigualdad
y la ciudadanía (\citeproc{ref-laffert_changes}{Laffert et~al., s.f.}).
El protocolo del estudio fue aprobado por un comité de ética, y se
registró la autorización de un adulto para la participación de cada
estudiante (\citeproc{ref-edumer_manual}{Proyecto Educación y
Meritocracia, s.f.}).

La encuesta es un panel de dos olas. La población objetivo son
estudiantes de escuelas municipales y particulares subvencionadas de las
regiones Metropolitana y de Valparaíso. Se encuestó a dos cohortes:
estudiantes de 6.º básico y 1.º medio en la primera ola, que fueron
seguidos en 7.º básico y 2.º medio en la segunda. Participaron nueve
escuelas (\citeproc{ref-laffert_changes}{Laffert et~al., s.f.}). La
muestra no es probabilística ni representativa por zona geográfica o
dependencia administrativa (\citeproc{ref-edumer_manual}{Proyecto
Educación y Meritocracia, s.f.}).

El cuestionario fue autoadministrado en línea (CAWI) y respondido por
los estudiantes en las salas de computación de sus escuelas, con una
duración aproximada de 20 minutos. La primera ola se aplicó entre
octubre de 2023 y mayo de 2024, y la segunda entre agosto y noviembre de
2024 (\citeproc{ref-edumer_manual}{Proyecto Educación y Meritocracia,
s.f.}). La primera ola incluye 902 estudiantes y la segunda 674. De
ellos, 639 participan en ambas olas, 263 solo en la primera y 35 solo en
la segunda, lo que implica una atrición de 29,2\%
(\citeproc{ref-edumer_manual}{Proyecto Educación y Meritocracia, s.f.}).

El análisis experimental (H4) utiliza los 902 estudiantes de la primera
ola, que es la única en que se aplicó el experimento. Las demás
hipótesis se contrastan con las observaciones de ambas olas, que suman
1.576 observaciones de 937 estudiantes.

La primera ola también encuestó a 180 apoderados y 26 docentes, cuyos
cuestionarios contienen el mismo ítem de igualdad de oportunidades
(\citeproc{ref-edumer_manual}{Proyecto Educación y Meritocracia, s.f.}).
Estas bases no forman parte del análisis principal, pero permiten un
análisis descriptivo complementario de la percepción de oportunidades en
la comunidad escolar.

\textbf{Experimento de encuesta.} El módulo experimental abre el
cuestionario de la primera ola, antes de todos los demás módulos. Los
estudiantes fueron asignados aleatoriamente, en proporción de 50\%
dentro de cada curso, a un grupo tratamiento (n = 444) o a un grupo
control (n = 458) (\texttt{asignacion})
(\citeproc{ref-edumer_manual}{Proyecto Educación y Meritocracia, s.f.}).
El grupo tratamiento leyó un texto titulado ``Tipos de colegio y acceso
a la universidad'', acompañado de una figura
(\citeproc{ref-castillo_experimento_2024}{Castillo, Carrasco, et~al.,
2024}):

\begin{quote}
Para entrar a la universidad se necesita tener buenas notas en el
colegio y también un buen puntaje en una prueba de selección de lenguaje
y matemáticas. En Chile existen colegios de distinto tipo: municipales,
subvencionados y privados. En los colegios municipales no se debe pagar
por estudiar, mientras que en los colegios privados las familias tienen
que pagar. Existen grandes diferencias en los puntajes de la prueba de
acceso a la universidad entre colegios municipales (no pagados) y
privados (pagados): 2 de cada 10 estudiantes de colegios municipales
logran entrar a la universidad, mientras en el caso de los colegios
privados son 8 de cada 10.
\end{quote}

El grupo control leyó un texto con la misma estructura, titulado ``Tipos
de colegio y acceso a deportes'', que presentaba la misma razón numérica
referida a la práctica deportiva: 2 de cada 10 estudiantes de colegios
municipales practican deportes, frente a 8 de cada 10 en colegios
privados (\citeproc{ref-castillo_experimento_2024}{Castillo, Carrasco,
et~al., 2024}). El diseño mantiene constantes el contraste entre tipos
de colegio y la magnitud de la diferencia, y varía solo el ámbito al que
se refiere: el acceso a la universidad o el deporte. Lo que se manipula,
por tanto, no es la información sobre la existencia de tipos de colegio
desiguales, sino que esa desigualdad se traduzca en oportunidades
educativas.

A continuación, cada grupo respondió una pregunta de verificación: ``¿De
qué tipo de colegio entran más estudiantes a la universidad?''
(\texttt{tratamiento}) o ``¿En qué tipo de colegio se practican más
deportes?'' (\texttt{control}). Al cierre del cuestionario se informó a
los estudiantes que las diferencias presentadas en el módulo
experimental habían sido exageradas
(\citeproc{ref-edumer_manual}{Proyecto Educación y Meritocracia, s.f.}).

\section{Operacionalización}\label{operacionalizaciuxf3n}

\subsection{Variables}\label{variables}

La variable dependiente es la percepción de igualdad de oportunidades,
medida con el ítem ``En Chile, todas las personas tienen las mismas
oportunidades para salir adelante'' (\texttt{p1\_9}), tomado de los
cuestionarios de la Agencia de Calidad de la Educación. Dado que se
trata de un ítem de cuatro categorías ordenadas, se analiza como
variable ordinal; como robustez, se trata también como escala numérica
de 1 a 4 y como variable dicotómica que distingue a quienes están de
acuerdo o muy de acuerdo de quienes están en desacuerdo o muy en
desacuerdo. El ítem contiguo del cuestionario, ``En Chile, todas las
personas obtienen lo que merecen'' (\texttt{p1\_10}), no se combina con
la variable dependiente, porque se refiere a la regla que vincula el
mérito con la recompensa y ha sido utilizado como indicador de
percepción de meritocracia
(\citeproc{ref-castillo_meritocracia_2019}{Castillo et~al., 2019}).

La percepción de meritocracia escolar (H1) se mide con tres ítems: ``En
esta escuela, quienes se esfuerzan obtienen buenas notas''
(\texttt{p2\_1}), ``En esta escuela, quienes son inteligentes obtienen
buenas notas'' (\texttt{p2\_2}) y ``En esta escuela, los/as estudiantes
obtienen las notas que merecen'' (\texttt{p2\_3}). Con ellos se
construye un índice que promedia las respuestas, de modo que valores más
altos indican una mayor percepción de que la escuela recompensa el
mérito. El rendimiento académico (H2) corresponde al promedio de notas
del año anterior (\texttt{p4}), en la escala chilena de 1,0 a 7,0. El
horizonte educativo (H3) se mide con la pregunta ``¿Cuál de los
siguientes niveles educacionales crees que vas a completar?''
(\texttt{p25}), dicotomizada entre quienes esperan completar una carrera
universitaria y quienes esperan completar la educación media o una
carrera en un instituto profesional o centro de formación técnica. La
información sobre el acceso a la universidad (H4) es la asignación
aleatoria al grupo tratamiento o control del experimento de la primera
ola (\texttt{asignacion}); la comprensión de la información se registra
con las preguntas de verificación (\texttt{tratamiento} y
\texttt{control}), considerando correcta la respuesta ``colegio
privado'' en ambos grupos. La trayectoria escolar (H5) se mide con la
ola de la encuesta, que distingue la primera de la segunda medición de
cada estudiante.

Las creencias meritocráticas (H6) se miden con la escala de percepciones
y preferencias sobre la meritocracia de Castillo et al.
(\citeproc{ref-castillo_multidimensional_2023}{2023}), que mide cada
dimensión con dos ítems (\citeproc{ref-laffert_changes}{Laffert et~al.,
s.f.}). La percepción de meritocracia (H6a) se mide con ``En Chile, las
personas son recompensadas por sus esfuerzos'' (\texttt{p1\_1}) y ``En
Chile, las personas son recompensadas por su inteligencia y habilidad''
(\texttt{p1\_2}); la percepción de privilegio (H6b), con ``En Chile, a
quienes tienen padres ricos les va mucho mejor en la vida''
(\texttt{p1\_3}) y ``En Chile, quienes tienen buenos contactos les va
mejor en la vida'' (\texttt{p1\_4}); la preferencia meritocrática (H6c),
con ``Quienes más se esfuerzan deberían obtener mayores recompensas que
quienes se esfuerzan menos'' (\texttt{p1\_5}) y ``Quienes poseen más
talento deberían obtener mayores recompensas que quienes poseen menos
talento'' (\texttt{p1\_6}); y la aceptación del privilegio (H6d), con
``Está bien que quienes tienen padres ricos les vaya bien en la vida''
(\texttt{p1\_7}) y ``Está bien que quienes tienen buenos contactos les
vaya bien en la vida'' (\texttt{p1\_8}). Cada dimensión se construye
como el promedio de sus dos ítems. La estructura de cuatro factores de
esta escala se ha validado en estos mismos datos y es invariante entre
olas dentro de los estudiantes, lo que permite interpretar su cambio en
el tiempo como sustantivo, aunque no entre cohortes, sobre todo en el
ítem de percepción del esfuerzo (\citeproc{ref-laffert_changes}{Laffert
et~al., s.f.}).

\subsection{Controles}\label{controles}

Los modelos incluyen un conjunto de variables de control orientadas a
aislar la relación entre los predictores y la percepción de igualdad de
oportunidades de diferencias en la composición de la muestra. El género
(\texttt{p20}) se incluye porque la percepción de meritocracia varía
según género en la población adulta
(\citeproc{ref-duru-bellat_whos_2012}{Duru-Bellat \& Tenret, 2012});
dado el reducido número de estudiantes que se identifica con otro
género, se distingue a los hombres de las mujeres y quienes se
identifican con otro género. El país de nacimiento del estudiante
(\texttt{p22\_1}), que distingue a quienes nacieron en Chile de quienes
nacieron en el extranjero, controla por una dimensión de desventaja
relevante en las escuelas estudiadas, en las que cerca de uno de cada
diez estudiantes nació fuera del país; como solo se preguntó en la
primera ola, se asigna a ambas observaciones de cada estudiante. La
cohorte (básica o media) controla por diferencias de edad y de nivel
escolar que podrían confundirse con el cambio entre olas, y por la falta
de invarianza de la escala de meritocracia entre cohortes
(\citeproc{ref-laffert_changes}{Laffert et~al., s.f.}). Finalmente, como
el experimento abre el cuestionario de la primera ola, todas las
respuestas de esa ola son posteriores a la asignación; por ello, los
modelos observacionales controlan por la asignación al tratamiento, que
toma el valor cero en la segunda ola.

El estudio no incluye una medida del nivel socioeconómico. El
cuestionario de estudiantes indaga la educación de la madre y del padre,
pero cerca de un cuarto de las observaciones no reporta la de la madre,
y la información de los apoderados solo puede vincularse a una fracción
minoritaria de los estudiantes. Se optó por no utilizar indicadores con
esa cobertura en lugar de introducir una medición poco confiable.

\subsection{Modelos (notaciones genneradas cn IA; pediennte de
confirmacion)}\label{modelos-notaciones-genneradas-cn-ia-pediennte-de-confirmacion}

Las hipótesis se contrastan con dos modelos de regresión logística
ordinal. El primero aprovecha el experimento de la primera ola para
estimar el efecto causal de la información sobre el acceso a la
universidad (H4). Para cada estudiante \(i\) y cada categoría \(k\) de
la variable dependiente:

\begin{equation}\protect\phantomsection\label{eq-modelo1}{
\text{logit}\,P(Y_i \le k) = \tau_k - \left(\beta_1 T_i + \gamma_{e[i]}\right)
}\end{equation}

donde \(Y_i\) es la percepción de igualdad de oportunidades, \(\tau_k\)
son los umbrales entre categorías, \(T_i\) es la asignación al
tratamiento y \(\gamma_{e[i]}\) son efectos fijos de escuela y cohorte,
los estratos en que se realizó la asignación. Como la asignación es
aleatoria, \(\beta_1\) estima el efecto de intención de tratar de la
información (H4: \(\beta_1 < 0\)), que se reporta también como efecto
marginal promedio sobre la probabilidad de estar de acuerdo con el ítem.

El segundo modelo utiliza las observaciones de ambas olas para
contrastar H1, H2, H3, H5 y H6. Como cada estudiante es observado hasta
dos veces, se estima un modelo multinivel de dos niveles, en que las
observaciones de cada ola (nivel 1) están anidadas en estudiantes (nivel
2), con un intercepto aleatorio por estudiante. Para el estudiante \(i\)
en la ola \(t\):

\begin{equation}\protect\phantomsection\label{eq-modelo2}{
\begin{aligned}
\text{logit}\,P(Y_{it} \le k) = \tau_k - \big(&\beta_1 \text{MeritEsc}_{it} + \beta_2 \text{Notas}_{it} + \beta_3 \text{Univ}_{it} + \beta_4 \text{Ola}_t \\
&+ \beta_5 \text{PercMerit}_{it} + \beta_6 \text{PercPriv}_{it} + \beta_7 \text{PrefMerit}_{it} + \beta_8 \text{AcepPriv}_{it} \\
&+ \boldsymbol{\delta}'\mathbf{X}_{it} + u_i\big)
\end{aligned}
}\end{equation}

donde \(\text{MeritEsc}_{it}\) es la percepción de meritocracia escolar,
\(\text{Notas}_{it}\) el promedio de notas, \(\text{Univ}_{it}\) la
expectativa de completar la universidad, \(\text{Ola}_t\) un indicador
de la segunda ola, \(\text{PercMerit}_{it}\), \(\text{PercPriv}_{it}\),
\(\text{PrefMerit}_{it}\) y \(\text{AcepPriv}_{it}\) las cuatro
dimensiones de la creencia meritocrática, \(\mathbf{X}_{it}\) el vector
de controles (género, país de nacimiento, cohorte y asignación al
tratamiento) y \(u_i \sim N(0, \sigma_u^2)\) el intercepto aleatorio del
estudiante, que recoge las diferencias estables entre estudiantes. Las
variables continuas se centran en su media general. Cada hipótesis
corresponde a un coeficiente: \(\beta_1 > 0\) (H1), \(\beta_2 > 0\)
(H2), \(\beta_3 > 0\) (H3), \(\beta_4 \neq 0\) (H5), cuyo signo permite
distinguir entre socialización (\(\beta_4 > 0\)) e ilustración
(\(\beta_4 < 0\)), \(\beta_5 > 0\) (H6a), \(\beta_6 < 0\) (H6b),
\(\beta_7 < 0\) (H6c) y \(\beta_8 > 0\) (H6d).

\section{Estrategia de análisis}\label{estrategia-de-anuxe1lisis}

El análisis se organiza en tres etapas.

\textbf{Descripción.} Se describirá la distribución de la percepción de
igualdad de oportunidades en cada ola y cohorte, y su correlación con la
percepción de meritocracia escolar y con las cuatro dimensiones de la
creencia meritocrática. Este paso permite evaluar empíricamente si los
estudiantes distinguen la igualdad de oportunidades de la meritocracia,
como sostiene el marco conceptual. De manera complementaria, se
comparará la distribución del ítem entre estudiantes, docentes y
apoderados de los mismos establecimientos.

\textbf{Análisis experimental (H4).} Se verificará el balance entre
ambos grupos en características que no pueden ser afectadas por la
información recibida, como género, cohorte, establecimiento y país de
nacimiento. El efecto de la información se estimará con el primer modelo
(Ecuación~\ref{eq-modelo1}), primero sin covariables y luego ajustando
por covariables. De manera exploratoria, se examinará si el efecto varía
según la cohorte y según la comprensión de la información.

\textbf{Modelo panel (H1, H2, H3, H5 y H6).} Las hipótesis
observacionales se contrastarán con el segundo modelo
(Ecuación~\ref{eq-modelo2}). Se reportarán los coeficientes como razones
de odds y la proporción de varianza explicada por los efectos fijos y
por el intercepto del estudiante (R² marginal y condicional). Para
distinguir la variación dentro de cada estudiante de las diferencias
estables entre estudiantes, los predictores que cambian entre olas se
descompondrán en su desviación respecto de la media del estudiante y en
dicha media. El componente dentro del estudiante indica si los cambios
de un mismo estudiante entre olas se asocian con cambios en su
percepción de igualdad de oportunidades, y no está afectado por
características estables no observadas.

\textbf{Robustez.} Se probarán modelos lineales sobre la escala de 1 a 4
y modelos logísticos sobre la versión dicotómica de la variable
dependiente; tambien se estimaran con difrerentes especificaciones de
los controles y covariables.

\bookmarksetup{startatroot}

\chapter{Resultados}\label{resultados}

\section{Análisis descriptivo}\label{anuxe1lisis-descriptivo}

\subsection{Percepción de igualdad de
oportunidades}\label{percepciuxf3n-de-igualdad-de-oportunidades}

Como la variable dependiente es un ítem ordinal de cuatro categorías, se
describe mediante su distribución completa y no mediante un promedio. La
Tabla~\ref{tbl-p19-ola} presenta la distribución por ola y la
Figura~\ref{fig-p19-distribucion} la desagrega por cohorte.

En la primera ola, el 39,3\% de los estudiantes está de acuerdo o muy de
acuerdo con que en Chile todas las personas tienen las mismas
oportunidades para salir adelante, mientras que el 60,7\% está en
desacuerdo o muy en desacuerdo. La categoría más frecuente es ``en
desacuerdo'' (36,2\%), que corresponde también a la mediana. La
percepción de igualdad de oportunidades es, por tanto, minoritaria entre
los estudiantes, aunque no marginal: cerca de cuatro de cada diez la
sostienen.

{

\begin{longtable}[]{@{}
  >{\raggedright\arraybackslash}p{(\linewidth - 4\tabcolsep) * \real{0.7714}}
  >{\raggedleft\arraybackslash}p{(\linewidth - 4\tabcolsep) * \real{0.1143}}
  >{\raggedleft\arraybackslash}p{(\linewidth - 4\tabcolsep) * \real{0.1143}}@{}}

\caption{\label{tbl-p19-ola}Percepción de igualdad de oportunidades
según ola. Frecuencias y porcentajes.}

\tabularnewline

\toprule\noalign{}
\begin{minipage}[b]{\linewidth}\raggedright
En Chile, todas las personas tienen las mismas oportunidades para salir
adelante
\end{minipage} & \begin{minipage}[b]{\linewidth}\raggedleft
Ola 1
\end{minipage} & \begin{minipage}[b]{\linewidth}\raggedleft
Ola 2
\end{minipage} \\
\midrule\noalign{}
\endhead
\bottomrule\noalign{}
\endlastfoot
Muy en desacuerdo & 220 (24,5\%) & 197 (29,2\%) \\
En desacuerdo & 325 (36,2\%) & 228 (33,8\%) \\
De acuerdo & 160 (17,8\%) & 141 (20,9\%) \\
Muy de acuerdo & 193 (21,5\%) & 108 (16,0\%) \\
Total & 898 & 674 \\

\end{longtable}

}

Las diferencias entre cohortes son marcadas. En la primera ola, el
52,9\% de los estudiantes de básica está de acuerdo o muy de acuerdo,
frente al 25,8\% de los de media. En básica, la categoría más frecuente
es ``muy de acuerdo'' (34,4\%), mientras que en media solo el 8,7\%
elige esa opción y el 74,2\% se ubica en el desacuerdo. La brecha se
mantiene en la segunda ola (48,3\% y 26,2\% de acuerdo,
respectivamente). Este patrón es consistente con el marco más crítico
que se ha observado entre los estudiantes de media al evaluar si el
esfuerzo es recompensado (\citeproc{ref-laffert_changes}{Laffert et~al.,
s.f.}), aunque la comparación entre cohortes debe interpretarse con
cautela, porque contrasta a grupos distintos de estudiantes y no el
cambio de los mismos estudiantes a lo largo del tiempo.

\begin{figure}

\caption{\label{fig-p19-distribucion}Percepción de igualdad de
oportunidades según ola y cohorte. Porcentajes; a la izquierda,
desacuerdo; a la derecha, acuerdo.}

\centering{

\pandocbounded{\includegraphics[keepaspectratio]{04-resultados_files/figure-pdf/fig-p19-distribucion-1.pdf}}

}

\end{figure}%

Entre olas, la percepción de igualdad de oportunidades tiende a
disminuir. El porcentaje de acuerdo pasa de 39,3\% a 36,9\%, y la
proporción que está muy en desacuerdo aumenta de 24,5\% a 29,2\%. Como
la composición de ambas olas difiere por la atrición, el cambio se
evalúa con los 635 estudiantes que responden en las dos mediciones
(Figura~\ref{fig-p19-transiciones}). De ellos, el 44,9\% mantiene la
misma respuesta, el 33,1\% se desplaza hacia el desacuerdo y el 22,0\%
hacia el acuerdo. La percepción es además relativamente inestable a
nivel individual: el 29,3\% de los estudiantes cambia de lado entre el
acuerdo y el desacuerdo de una ola a otra. Este descenso es consistente
con la hipótesis de ilustración (H5), que se evalúa formalmente en los
modelos multivariados.

\begin{figure}

\caption{\label{fig-p19-transiciones}Transiciones en la percepción de
igualdad de oportunidades entre la ola 1 y la ola 2. Estudiantes que
responden ambas olas; porcentajes sobre el total de esos estudiantes.
Los flujos tienen el color de la respuesta en la ola 1.}

\centering{

\pandocbounded{\includegraphics[keepaspectratio]{04-resultados_files/figure-pdf/fig-p19-transiciones-1.pdf}}

}

\end{figure}%

\bookmarksetup{startatroot}

\chapter{Discusión}\label{discusiuxf3n}

Pendiente.

\section{Limitaciones}\label{limitaciones}

Por lo pronto no hay forma consistente de operacionalizar el NSE, poca
varianza en educación de los padres, discordancia entre los niveles
educativos reportados por los estudiantes y los apoderados, cantidad de
libros como proxy de capital cultural tampoco en este caso reporta
efectos significativos en ninguna de sus especificaciones, ni en versión
ordinal, ni en más de veinticinco versus el resto, los factores de
meritocracia estarían entrando como índices ponderados; se están sumando
los errores, habría que considerar algún tipo de análisis de robustez,
me da la impresión de que la inclusión del experimento está interesante
más no del todo justificada, de cualquier forma sirve para controlar al
menos en la ola uno si la info que se entrega a los estudiantes tiene
algún efecto sobre el resto de respuestas de la encuesta; ``¿Conocer
cifras de desigualdad en el acceso a la universidad afecta la percepcion
de igualdad de oportunidades?'' se me ocurre que tal vez podría
imputarse el control a la ola dos, si los estudiantes en algún momento
se les entregó la información sobre desigualdad, podría ser que para el
cuestionario dos se hayan acordado y también tenga efecto sobre sus
respuestas.

\bookmarksetup{startatroot}

\chapter{Conclusiones}\label{conclusiones}

Síntesis y proyecciones futuras.

\bookmarksetup{startatroot}

\chapter{Referencias}\label{referencias}

\protect\phantomsection\label{refs}
\begin{CSLReferences}{1}{0}
\bibitem[\citeproctext]{ref-allouch_ethique_2016}
Allouch, A., Brown, P., Power, S., \& Tholen, G. (2016). L'éthique des
élites scolaires. Du mérite à la responsabilité chez les étudiants de
Sciences Po Paris et de l'Université d'Oxford. \emph{L'Année
sociologique}, \emph{66}(1), 193-224.
\url{https://doi.org/10.3917/anso.161.0193}

\bibitem[\citeproctext]{ref-alwin_political_1991}
Alwin, D. F., Cohen, R. L., \& Newcomb, T. M. (1991). \emph{Political
Attitudes over the Life Span: The Bennington Women after Fifty Years}.
University of Wisconsin Press.

\bibitem[\citeproctext]{ref-ashley_understanding_2017}
Ashley, L., \& Empson, L. (2017). Understanding social exclusion in
elite professional service firms: field level dynamics and the
{«professional project»}. \emph{Work, Employment and Society},
\emph{31}(2), 211-229. \url{https://doi.org/10.1177/0950017015621480}

\bibitem[\citeproctext]{ref-baer_education_1982}
Baer, D. E., \& Lambert, R. D. (1982). Education and Support for
Dominant Ideology. \emph{Canadian Review of Sociology / Revue canadienne
de sociologie}, \emph{19}(2), 173-195.
\url{https://doi.org/10.1111/j.1755-618X.1982.tb00859.x}

\bibitem[\citeproctext]{ref-bathmaker_degree_2016}
Bathmaker, A.-M., Ingram, N., Abrahams, J., Hoare, A., Waller, R., \&
Bradley, H. (2016). \emph{Higher Education, Social Class and Social
Mobility: The Degree Generation}. Palgrave Macmillan.

\bibitem[\citeproctext]{ref-batruch_belief_2023}
Batruch, A., Jetten, J., Van De Werfhorst, H., Darnon, C., \& Butera, F.
(2023). Belief in {School Meritocracy} and the {Legitimization} of
{Social} and {Income Inequality}. \emph{Social Psychological and
Personality Science}, \emph{14}(5), 621-635.
\url{https://doi.org/10.1177/19485506221111017}

\bibitem[\citeproctext]{ref-binder_symbolically_2019}
Binder, A. J., \& Abel, A. R. (2019). Symbolically Maintained
Inequality: How Harvard and Stanford Students Construct Boundaries among
Elite Universities. \emph{Sociology of Education}, \emph{92}(1), 41-58.
\url{https://doi.org/10.1177/0038040718821073}

\bibitem[\citeproctext]{ref-bourdieu_reproduction_1990}
Bourdieu, P., \& Passeron, J.-C. (1990). \emph{Reproduction in
Education, Society and Culture} (2.ª ed.). Sage.

\bibitem[\citeproctext]{ref-brunner_lucha_2012}
Brunner, J. J. (2012). La lucha por la educación de las elites: campo y
canales formativos. \emph{Revista UDP}, \emph{9}, 119-143.

\bibitem[\citeproctext]{ref-brunner_dinamicas_2016}
Brunner, J. J., \& Ganga, F. A. (2016). Dinámicas de transformación en
la educación superior latinoamericana: desafíos para la gobernanza.
\emph{Opción}, \emph{32}(80), 12-35.

\bibitem[\citeproctext]{ref-castillo_experimento_2024}
Castillo, J. C., Carrasco, K., Nuñez, M. F., \& Laffert, A. (2024).
\emph{Meritocracia en la etapa escolar: un experimento de encuesta}.
Presentación en el Congreso Chileno de Sociología, Santiago, 2-4 de
mayo. \url{https://github.com/educacion-meritocracia/experimento}

\bibitem[\citeproctext]{ref-castillo_multidimensional_2023}
Castillo, J. C., Iturra, J., Maldonado, L., Atria, J., \& Meneses, F.
(2023). A {Multidimensional Approach} for {Measuring Meritocratic
Beliefs}: {Advantages}, {Limitations} and {Alternatives} to the {ISSP
Social Inequality Survey}. \emph{International Journal of Sociology},
\emph{53}(6), 448-472.
\url{https://doi.org/10.1080/00207659.2023.2274712}

\bibitem[\citeproctext]{ref-castillo_percepcion_2012}
Castillo, J. C., Miranda, D., \& Carrasco, D. (2012). Percepción de
desigualdad económica en Chile: medición, diferencias y determinantes.
\emph{Psykhe}, \emph{21}(1), 99-114.

\bibitem[\citeproctext]{ref-castillo_socialization_2024}
Castillo, J. C., Salgado, M., Carrasco, K., \& Laffert, A. (2024). The
{Socialization} of {Meritocracy} and {Market Justice Preferences} at
{School}. \emph{Societies}, \emph{14}(11), 214.
\url{https://doi.org/10.3390/soc14110214}

\bibitem[\citeproctext]{ref-castillo_meritocracia_2019}
Castillo, J. C., Torres, A., Atria, J., \& Maldonado, L. (2019).
Meritocracia y Desigualdad Económica: {Percepciones}, Preferencias e
Implicancias. \emph{Revista Internacional de Sociología}, \emph{77}(1),
e117. \url{https://doi.org/10.3989/ris.2019.77.1.17.114}

\bibitem[\citeproctext]{ref-darnon_where_2018}
Darnon, C., Wiederkehr, V., Dompnier, B., \& Martinot, D. (2018).
{«{Where} There Is a Will, There Is a Way»}: {Belief} in School
Meritocracy and the Social-class Achievement Gap. \emph{British Journal
of Social Psychology}, \emph{57}(1), 250-262.
\url{https://doi.org/10.1111/bjso.12214}

\bibitem[\citeproctext]{ref-dickinson_adolescent_1990}
Dickinson, J. (1990). Adolescent Representations of Socio-economic
Status. \emph{British Journal of Developmental Psychology}, \emph{8}(4),
351-371. \url{https://doi.org/10.1111/j.2044-835X.1990.tb00850.x}

\bibitem[\citeproctext]{ref-duru-bellat_whos_2012}
Duru-Bellat, M., \& Tenret, E. (2012). Who's for Meritocracy?
{Individual} and Contextual Variations in the Faith. \emph{Comparative
Education Review}, \emph{56}(2), 223-247.
\url{https://doi.org/10.1086/661290}

\bibitem[\citeproctext]{ref-goudeau_hidden_2017}
Goudeau, S., \& Croizet, J.-C. (2017). Hidden Advantages and
Disadvantages of Social Class: {How} Classroom Settings Reproduce Social
Inequality by Staging Unfair Comparison. \emph{Psychological Science},
\emph{28}(2), 162-170. \url{https://doi.org/10.1177/0956797616676600}

\bibitem[\citeproctext]{ref-laffert_changes}
Laffert, A., Castillo, J. C., Canales, R., Urzúa, T., \& Carrasco, K.
(s.f.). \emph{Changes in meritocratic beliefs among school-age students
over time and across cohorts}.

\bibitem[\citeproctext]{ref-mijs_stratified_2016}
Mijs, J. J. B. (2016a). Stratified Failure: {Educational} Stratification
and Students' Attributions of Their Mathematics Performance in 24
Countries. \emph{Sociology of Education}, \emph{89}(2), 137-153.
\url{https://doi.org/10.1177/0038040716636434}

\bibitem[\citeproctext]{ref-mijs2016}
Mijs, J. J. B. (2016b). \emph{The Unfulfillable Promise of Meritocracy:
Three Lessons and Their Implications for Justice in Education}.

\bibitem[\citeproctext]{ref-munoz_taxonomia_2013}
Muñoz, M., \& Blanco, C. (2013). Una taxonomía de las universidades
chilenas. \emph{Calidad en la Educación}, \emph{38}.

\bibitem[\citeproctext]{ref-edumer_manual}
Proyecto Educación y Meritocracia. (s.f.). \emph{Encuesta Panel
Educación y Meritocracia (EDUMER): manual metodológico}.
\url{https://educacion-meritocracia.github.io/encuesta-edumer/}

\bibitem[\citeproctext]{ref-resh2010}
Resh, N., \& Sabbagh, C. (2010). Sense of Justice about Grades in
School: Is It Stratified like Academic Achievement? \emph{Social
Psychology of Education}.

\bibitem[\citeproctext]{ref-villalobos_percepciones_2022}
Villalobos, C., Kuzmanic, D., Valenzuela, J. P., \& Quaresma, M. L.
(2022). Percepciones sociales de los estudiantes de educación superior
de élite: un estudio del caso chileno. \emph{Psykhe}, \emph{31}(2),
1-18. \url{https://doi.org/10.7764/psykhe.2020.22549}

\bibitem[\citeproctext]{ref-young1958}
Young, M. (1958). \emph{The Rise of the Meritocracy, 1870--2033}.
Thames; Hudson.

\end{CSLReferences}

\cleardoublepage
\phantomsection
\addcontentsline{toc}{part}{Apéndices}
\appendix

\chapter{Instrumentos de
levantamiento}\label{instrumentos-de-levantamiento}

Incluye cuestionarios, guias de entrevista, etc.

\chapter{Tablas adicionales}\label{tablas-adicionales}

Tablas extendidas, tests de robustez, etc.


\backmatter

% --- Cuerpo del libro (capítulos) en arábigos ---
\mainmatter
\pagestyle{scrheadings}

% --- (Opcional) Estilo del índice general ---
% \addtocontents{toc}{\protect\thispagestyle{plain}}


\end{document}
